% ================================================================
\chapter{Ordenação em tempo linear}\label{cap:sortlin}\label{sec:sortlin}
% ================================================================
\begin{objetivos}
\begin{itemize}
  \item Como escapar do limite $\Omega(n\log n)$ usando a estrutura das chaves.
  \item Counting sort, radix sort, bucket sort e American Flag sort.
  \item O papel da estabilidade no radix sort.
  \item Um guia para escolher o algoritmo de ordenação adequado a cada situação.
\end{itemize}
\end{objetivos}

O limite $\Omega(n\log n)$ do capítulo anterior vale para algoritmos que só \emph{comparam} elementos. Qualquer algoritmo de ordenação em tempo linear precisa, portanto, fazer no máximo $O(n)$ comparações e tomar suas decisões de outra forma. Se soubermos algo sobre as chaves --- que são inteiros num intervalo pequeno, números com poucos dígitos ou reais uniformemente distribuídos ---, podemos usá-las diretamente como \emph{índices} e ordenar em tempo linear. O preço é memória extra e hipóteses sobre os dados.

\section{Counting sort}
O counting sort ordena inteiros no intervalo $[0,k]$. Ele conta quantas vezes cada valor ocorre, acumula as contagens para saber quantos elementos são menores ou iguais a cada valor e, com isso, coloca cada elemento diretamente na sua posição final.
\begin{pseudo}{Counting-Sort$(A,n,k)$}
\begin{algorithmic}
\For{$i:=0$ \textbf{até} $k$} \State $C[i]:=0$ \EndFor
\For{$j:=1$ \textbf{até} $n$} \State $C[A[j]]:=C[A[j]]+1$ \EndFor \Comment{$C[i]=\#\{j: A[j]=i\}$}
\For{$i:=1$ \textbf{até} $k$} \State $C[i]:=C[i]+C[i-1]$ \EndFor \Comment{$C[i]=\#\{j: A[j]\le i\}$}
\For{$j:=n$ \textbf{decrescendo até} $1$} \Comment{de trás para frente: estável}
  \State $B[C[A[j]]]:=A[j]$;\ $C[A[j]]:=C[A[j]]-1$
\EndFor
\State \Return $B$
\end{algorithmic}
\end{pseudo}
O tempo e a memória são $\Theta(n+k)$, o que é \textbf{linear quando $k=O(n)$}. Quando $k$ é muito maior que $n$ (por exemplo, $k=\Theta(2^n)$), o algoritmo fica exponencial no tamanho da entrada: é \textbf{pseudopolinomial}. Ele é \textbf{estável}, porque, percorrendo a entrada de trás para frente, a última ocorrência de cada valor vai para a última posição do seu bloco; não é in-place, pois usa os vetores $B$ e $C$; e não faz nenhuma comparação entre elementos.

\section{Radix sort}
O radix sort ordena inteiros com $d$ dígitos na base $b$ \textbf{dígito a dígito, do menos significativo para o mais significativo}, usando em cada passada um método de ordenação \textbf{estável}. A versão abaixo usa uma fila para cada valor de dígito.
\begin{pseudo}{Radix sort com filas}
\begin{algorithmic}
\For{$i:=1$ \textbf{até} $d$}
  \For{$j:=1$ \textbf{até} $n$}
    \State $k:=$ $i$-ésimo dígito menos significativo de $L[j]$ na base $b$;\ insira $L[j]$ no fim da fila $F[k]$
  \EndFor
  \State $j:=1$
  \For{$k:=0$ \textbf{até} $b-1$}
    \While{$F[k]$ não está vazia} \State retire o início de $F[k]$ e coloque-o em $L[j]$;\ $j:=j+1$ \EndWhile
  \EndFor
\EndFor
\end{algorithmic}
\end{pseudo}
\textbf{Por que a estabilidade é necessária?} Depois da passada pelo dígito $i$, elementos com o mesmo dígito $i$ devem manter entre si a ordem obtida pelos dígitos $1,\dots,i-1$. É isso que garante, por indução, que após a passada $i$ a lista está ordenada pelos $i$ últimos dígitos. As filas preservam a ordem de chegada naturalmente; outra opção é usar o counting sort em cada dígito.

O tempo é $\Theta(d(n+b))$: com $b$ e $d$ constantes, $\Theta(n)$. Uma escolha útil é usar a própria base $b=n$: inteiros menores que $n^c$ têm $c$ dígitos nessa base, e o radix sort os ordena em $\Theta(cn)$.

\begin{exemplo}[Radix sort na base 10]
Ordenemos $102,313,690,131,222,670,690,130,079,102$.\\
Pelas unidades: $690\ 670\ 690\ 130\ 131\ 102\ 222\ 102\ 313\ 079$;\\
pelas dezenas: $102\ 102\ 313\ 222\ 130\ 131\ 670\ 079\ 690\ 690$;\\
pelas centenas: $079\ 102\ 102\ 130\ 131\ 222\ 313\ 670\ 690\ 690$.
\end{exemplo}

\paragraph{Chaves de tamanhos variados.} Se os números têm quantidades diferentes de dígitos, $d_1,\dots,d_n$, completar todos com zeros até o maior tamanho desperdiça trabalho. É possível ordenar em tempo $O(\sum_id_i)$, linear no tamanho da entrada, agrupando os números pelo número de dígitos (um número com menos dígitos é menor) e fazendo cada número participar apenas das passadas correspondentes aos seus próprios dígitos. Para \textbf{strings}, ordena-se da direita para a esquerda completando as menores com um símbolo ``menor que todos'', ou usa-se a versão que começa pelo caractere mais significativo, descrita a seguir.

\section{Bucket sort}
O bucket sort supõe que as chaves são reais \textbf{uniformemente distribuídos} em $[0,1)$. Divide-se o intervalo em $r$ baldes de mesmo tamanho, coloca-se $A[i]$ no balde $\floor{r\,A[i]}$, ordena-se cada balde (por exemplo, por inserção) e concatenam-se os baldes. Com entrada uniforme, cada balde recebe em média $n/r$ elementos; com $r=\Theta(n)$, o tempo \textbf{esperado} é $\Theta(n)$, porque, sendo $n_i$ o número de elementos no balde $i$, vale $E[n_i^2]=2-1/n$, e $\sum_iE[O(n_i^2)]=O(n)$. O pior caso, quando todos os elementos caem no mesmo balde, é o do algoritmo usado dentro dos baldes.

\section{American Flag sort}
O American Flag sort é um radix sort que começa pelo dígito \textbf{mais} significativo e trabalha \textbf{in-place}. Ele conta quantos elementos há com cada valor do dígito (um vetor $C$ de tamanho $b$), calcula onde começa o bloco de cada valor e \emph{permuta no próprio vetor} cada elemento para o seu bloco, por meio de trocas; depois, chama a si mesmo recursivamente em cada bloco, com o dígito seguinte. A memória auxiliar se resume aos contadores: $O(b)$ por nível de recursão, $O(b\cdot d)$ no total, independente de $n$. As trocas mudam a ordem relativa de elementos iguais, de modo que o algoritmo \textbf{não é estável} --- e não precisa ser, porque cada bloco é ordenado pelos dígitos seguintes de forma independente.

% ------------------------------------------------------------------
\section{Qual algoritmo de ordenação usar?}\label{sec:escolhasort}
\begin{center}
\small
\begin{tabularx}{\linewidth}{@{}lX@{}}
\toprule
Situação & Escolha\\
\midrule
Inteiros num intervalo pequeno ($k=O(n)$) & \textbf{Counting sort}, $\Theta(n+k)$\\
Inteiros ou identificadores com poucos dígitos ($d$ fixo) e $k$ enorme & \textbf{Radix sort}, $\Theta(d(n+b))$\\
Reais uniformemente distribuídos num intervalo & \textbf{Bucket sort}, $\Theta(n)$ esperado\\
Chaves arbitrárias, com garantia de tempo e/ou estabilidade & \textbf{Mergesort}, $\Theta(n\log n)$, memória $\Theta(n)$\\
Chaves arbitrárias, pouca memória, com garantia de tempo & \textbf{Heapsort}, $\Theta(n\log n)$, in-place\\
Chaves arbitrárias, desempenho médio na prática & \textbf{Quicksort} com pivô aleatório ou mediana de três\\
Vetor quase ordenado, ou $n$ pequeno & \textbf{Inserção}\\
Dados em lista encadeada ou em arquivos externos & \textbf{Mergesort}\\
\bottomrule
\end{tabularx}
\end{center}

\begin{exemplo}[Escolhendo o algoritmo]
(A) $10^7$ notas inteiras entre 0 e 100: \textbf{counting sort}, com 101 contadores e tempo $\Theta(n)$; uma ordenação por comparação faria cerca de $n\log_2n\approx2{,}3\cdot10^8$ comparações. (B) $10^6$ números reais uniformes em $[0,1)$: \textbf{bucket sort}, $\Theta(n)$ esperado; counting e radix não se aplicam a reais. (C) $10^5$ registros com chaves arbitrárias: só resta comparar, e o \textbf{mergesort} dá garantia de $\Theta(n\log n)$ e estabilidade.
\end{exemplo}

\begin{chave}
Counting sort: $\Theta(n+k)$, estável (último laço de trás para frente), linear se $k=O(n)$. Radix sort (do menos para o mais significativo): $\Theta(d(n+b))$, exige passadas estáveis. Bucket sort: $\Theta(n)$ esperado para entrada uniforme. American Flag sort: radix pelo dígito mais significativo, in-place e não estável. Para chaves arbitrárias, o limite $\Omega(n\log n)$ continua valendo.
\end{chave}

% ----------------------------------------------------------------
\section{Exercícios}
\excurso
\begin{exercicio}{sl-t2}[Avaliação]
Uma empresa precisa ordenar grandes volumes de dados com características diferentes. Para cada cenário, escolha o algoritmo mais adequado entre counting sort, radix sort, bucket sort, mergesort e quicksort, justificando por tempo, espaço e características dos dados: (A) $10^7$ notas entre 0 e 100; (B) $10^6$ identificadores numéricos de 12 dígitos; (C) $10^5$ números reais uniformemente distribuídos entre 0 e 1; (D) $10^6$ inteiros com valores entre $-10^9$ e $10^9$; (E) $10^5$ registros com chaves arbitrárias e sem restrição de domínio.
\end{exercicio}
\begin{exercicio}{sl-var}[Slides]
Modifique o radix sort para ordenar inteiros com quantidades diferentes de dígitos em tempo $O\big(\sum_{i=1}^nd_i\big)$, onde $d_i$ é o número de dígitos do $i$-ésimo inteiro. Faça o mesmo para o American Flag sort.
\end{exercicio}
\begin{exercicio}{sl-str}[Slides]
Que modificações são necessárias no radix sort e no American Flag sort para ordenar \emph{strings} em vez de inteiros?
\end{exercicio}

\excompl
\begin{exercicio}{sl-count}[Complementar]
Execute o counting sort em $A=[2,5,3,0,2,3,0,3]$ com $k=5$: mostre o vetor $C$ após a contagem e após as somas acumuladas, e a saída $B$.
\end{exercicio}
\begin{exercicio}{sl-radix}[Complementar]
Execute o radix sort na base 10 sobre $329,457,657,839,436,720,355$, mostrando a lista após cada passada.
\end{exercicio}
\begin{exercicio}{sl-n3}[Complementar]
Mostre como ordenar $n$ inteiros do intervalo $[0,n^3-1]$ em tempo $O(n)$.
\end{exercicio}
\begin{exercicio}{sl-frente}[Complementar]
No último laço do counting sort, o que muda se $j$ for de $1$ até $n$ (em vez de $n$ até $1$)? O resultado continua ordenado? E o radix sort que usa esse counting sort em cada passada?
\end{exercicio}
\begin{exercicio}{sl-bucket}[Complementar]
Qual o pior caso do bucket sort com ordenação por inserção nos baldes? Como garantir $O(n\log n)$ no pior caso sem perder o tempo esperado $\Theta(n)$?
\end{exercicio}
\begin{exercicio}{sl-msd}[Complementar]
Por que o radix sort que começa pelo dígito menos significativo precisa de passadas estáveis, enquanto o American Flag sort, que começa pelo mais significativo, funciona mesmo não sendo estável?
\end{exercicio}
\begin{exercicio}{sl-neg}[Complementar]
Como usar o counting sort para ordenar inteiros no intervalo $[-m,m]$? Qual o custo?
\end{exercicio}
\begin{exercicio}{sl-n2}[Complementar]
Para ordenar $n$ inteiros no intervalo $[0,n^2-1]$, compare o counting sort direto com o radix sort na base $n$.
\end{exercicio}
\begin{exercicio}{sl-strtrace}[Complementar]
Ordene pelo radix sort as palavras COW, DOG, SEA, RUG, ROW, MOB, BOX, TAB, BAR, EAR, TAR, DIG, BIG, TEA, NOW, FOX, mostrando a lista após cada uma das três passadas.
\end{exercicio}
\begin{exercicio}{sl-buckettrace}[Complementar]
Execute o bucket sort com 10 baldes sobre $0{,}78;\ 0{,}17;\ 0{,}39;\ 0{,}26;\ 0{,}72;\ 0{,}94;\ 0{,}21;\ 0{,}12;\ 0{,}23;\ 0{,}68$. Mostre o conteúdo de cada balde não vazio.
\end{exercicio}
\begin{exercicio}{sl-base}[Complementar]
Para ordenar $10^6$ inteiros de 32 bits com o radix sort, compare as bases $2^4$, $2^8$ e $2^{16}$: quantas passadas cada uma faz e qual o custo aproximado de cada passada? Qual parece a melhor escolha?
\end{exercicio}
\begin{exercicio}{sl-prova}[Complementar]
Prove por indução no número de passadas que o radix sort (do dígito menos significativo para o mais significativo, com passadas estáveis) ordena corretamente.
\end{exercicio}
\begin{exercicio}{sl-registros}[Complementar]
Uma lista de alunos já está em ordem alfabética e deve ser ordenada pela nota (um inteiro de 0 a 10), mantendo em ordem alfabética os alunos com a mesma nota. Que algoritmo usar e por quê?
\end{exercicio}
\begin{exercicio}{sl-datas}[Complementar]
Datas no formato (dia, mês, ano) devem ser ordenadas cronologicamente com o radix sort. Em que ordem devem ser feitas as passadas e qual a base de cada uma?
\end{exercicio}
\begin{exercicio}{sl-dup}[Complementar]
Dados $n$ inteiros no intervalo $[0,n^2]$, decida em tempo $O(n)$ se há elementos repetidos.
\end{exercicio}
\begin{exercicio}{sl-intervalo}[Complementar]
Dados $n$ inteiros no intervalo $[0,k]$, faça um pré-processamento em tempo $O(n+k)$ que permita responder em $O(1)$ a perguntas do tipo ``quantos elementos estão no intervalo $[a,b]$?''.
\end{exercicio}
\begin{exercicio}{sl-bucketruim}[Complementar]
Dê uma entrada de $n$ números em $[0,1)$ para a qual o bucket sort com $n$ baldes e inserção nos baldes gasta $\Theta(n^2)$.
\end{exercicio}
\begin{exercicio}{sl-flag}[Complementar]
Quanta memória auxiliar o American Flag sort usa para ordenar $n$ inteiros de $d$ dígitos na base $b$? Por que se diz que ele é in-place, se usa vetores de contadores?
\end{exercicio}
